\section{Introduction}
\label{sec:introduction}

\begin{layman}
\textbf{What is this document?}\par

In late 2024, the mathematician Kevin Buzzard issued a public
challenge: formalize, in the proof assistant Lean, a small but
pivotal piece of modern mathematics — the construction and basic
properties of the \emph{Jacobian} of a compact Riemann surface. The
Jacobian is one of the most studied objects in nineteenth- and
twentieth-century geometry. To a curved two-dimensional surface
like a doughnut or a pretzel, it attaches a higher-dimensional
companion that records all the surface's topological information in
a clean, linear-algebraic form. From genus-one curves and elliptic
integrals to the Birch--Swinnerton-Dyer conjecture and Wiles's proof
of Fermat's Last Theorem, every road in modern arithmetic geometry
passes through the Jacobian.

The challenge has a peculiar structure. Its public goal is a short
Lean file, \code{Jacobian/Challenge.lean}, listing perhaps two dozen
named definitions and theorems. Filling those entries with literal
\code{sorry} placeholders takes about thirty lines. The hard part
is that four of the theorems are deliberately constructed to defeat
any cheat — any attempt to define \code{Jacobian X := PUnit} (a
single point) or \code{genus X := 0} (always zero) and call it done.
They force the genus to track topology, the Jacobian to be a
genuine compact complex Lie group, the Abel--Jacobi map to be
injective, and pullback and pushforward to interact through the
classical trace identity. To meet them one must build real
mathematical infrastructure: not just statements but the
several-layer stack of definitions, lemmas, and analytic background
that make the statements provable.

This document is the \emph{informal proof roadmap} for that
formalization. It plays two roles at once. As a piece of
mathematical exposition, it lays out a complete paper-style proof
that the classical analytic construction satisfies the challenge:
every \code{sorry} in \code{Challenge.lean} corresponds to an
English theorem here, with sources and dependencies attached. As a
planning document, it gives every theorem a stable label and,
where the Lean has been written, a pointer to the corresponding
declaration. The dependency graph between labels is what drives the
project's automated proof work; a remote tool called
\emph{Aristotle} picks bounded leaves off the graph and discharges
them one at a time.

The construction we use is one of several available routes — the
\emph{analytic period-lattice} construction. Pick a compact connected
Riemann surface \(X\). Form the finite-dimensional complex vector
space \(H^0(X,\Omega^1_X)\) of holomorphic 1-forms on \(X\). Take its
linear dual. Inside that dual, find the \emph{period lattice}: the
discrete grid of vectors obtained by integrating each form along
each integer-homology loop. Quotient by this lattice. The result, a
compact complex torus of dimension \(g(X)\), is the Jacobian. Every
challenge property — group structure, manifold structure, every
anti-hack theorem — falls out of this single construction once
enough analytic background has been established.

The nine-section spine of the document mirrors that construction.
Section~1 sets up the algebraic machinery (vanishing orders,
divisors) shared with the algebraic-geometric route. Section~2
develops \(H^0(X,\Omega^1_X)\), its finite-dimensionality,
Riemann--Roch, and the genus-zero classification. Section~3 develops
the \emph{polygonal model} — the classical unfolding of the surface
onto a flat \(4g\)-gon — that underlies all later integration
arguments. Section~4 builds the period pairing and the Riemann
bilinear relations. Section~5 constructs the Jacobian via the
complex-torus package. Section~6 establishes the Abel--Jacobi map
and its injectivity. Section~7 develops degree, trace, pullback and
pushforward, and the trace identity. Section~8 assembles the public
API. Section~9 maps each Lean declaration in \code{Challenge.lean}
back to its informal source.

A reader can approach the document at three levels. Skim only the
section headings for a high-level sense of the construction. Read
the boldfaced theorem statements and surrounding prose for a
working mathematician's tour. Read the \textbf{Plain English}
insets — like this one — for a general-audience narrative version,
free of any prerequisite beyond curiosity. The formal-only PDF
strips these insets; the companion build keeps them.
\end{layman}

The public target is the API in \code{Jacobian/Challenge.lean}. This
document gives an informal proof that the API is satisfied by the classical
analytic construction of the Jacobian of a compact connected Riemann surface.

The construction is:
\[
  \Jac(X)=H^{0}(X,\Omega^{1}_{X})^{\vee}/\Pi_X(H_{1}(X,\Z)),
\]
where \(\Pi_X(\gamma)(\omega)=\int_{\gamma}\omega\). The main mathematical
work is to prove that the period image is a full lattice, to relate the genus
to the dimension of holomorphic one-forms, and to prove the Abel-Jacobi and
trace identities needed by the challenge.

The document has two roles. First, it is a readable classical proof. Second, it
is a naming and dependency plan for formalization: every major theorem has a
stable label and, when possible, a \code{Lean target} annotation matching the
current repository.
